% ============================================================
%  A6 - Großes Projekt: Hauptdatei
%  Kompilieren: latexmk -pdf main.tex   (oder mehrmals pdflatex)
%  Struktur:
%    main.tex            Präambel und Reihenfolge der Kapitel
%    abkuerzungen.tex    Definition der Abkürzungen
%    kapitel/*.tex       je Kapitel eine Datei
% ============================================================
\documentclass[a4paper,11pt]{report}
\usepackage[T1]{fontenc}
\usepackage[ngerman]{babel}
\usepackage{graphicx}
\usepackage{booktabs}

% --- Anmerkungen im Rand ------------------------------------
% \todo{...} setzt eine Notiz in den Rand, \listoftodos listet alle auf.
% Für die Abgabe: Option "disable" ergänzen, dann verschwinden alle Notizen.
\usepackage[textsize=footnotesize]{todonotes}
\setlength{\marginparwidth}{2.5cm}

% --- Reihenfolge ist wichtig: hyperref -> glossaries -> cleveref ---
\usepackage[hidelinks]{hyperref}     % klickbare Verweise

% Abkürzungsverzeichnis
%   acronym: eigene Liste für Abkürzungen
%   nomain:  kein allgemeines Glossar
% \makenoidxglossaries sortiert innerhalb von LaTeX. Es ist kein
% Zusatzprogramm nötig, daher läuft es in Overleaf und lokal gleich.
\usepackage[acronym,nomain]{glossaries}
\makenoidxglossaries
% Abkürzungen: \newacronym{Schlüssel}{Kurzform}{Langform}
% Im Text: \gls{ci}   -> beim ersten Mal "Continuous Integration (CI)", danach "CI"
%          \acrshort{ci} -> immer "CI",  \acrlong{ci} -> immer die Langform
\newacronym{ci}{CI}{Continuous Integration}
\newacronym{cd}{CD}{Continuous Delivery}
\newacronym{pdf}{PDF}{Portable Document Format}


% Verweise mit automatischem Namen: \cref{fig:x} ergibt "Abbildung 1"
\usepackage[ngerman]{cleveref}

% --- Nur einzelne Kapitel übersetzen ------------------------
% Spart Zeit bei großen Dokumenten. Seitenzahlen und Verweise der
% ausgelassenen Kapitel bleiben erhalten (aus den .aux-Dateien).
% \includeonly{kapitel/grundlagen}

\title{Projektbericht}
\author{Max Mustermann}
\date{\today}

\begin{document}

\maketitle
\tableofcontents
\listoftodos

% \include beginnt immer eine neue Seite und legt pro Datei eine
% eigene .aux-Datei an. Dateiname ohne Endung angeben.
\chapter{Einleitung}
\label{ch:einleitung}

Dieses Dokument zeigt, wie ein größeres \LaTeX-Projekt aufgebaut wird.
Die Grundlagen folgen in \cref{ch:grundlagen}, das Fazit in \cref{ch:fazit}.

\todo{Motivation ergänzen}
Bei jedem Push wird das Dokument per \gls{ci} neu gebaut.
Das Ergebnis ist eine \gls{pdf}-Datei.

\chapter{Grundlagen}
\label{ch:grundlagen}

\section{Verweise mit cleveref}
\label{sec:verweise}

\Cref{tab:befehle} fasst die wichtigsten Befehle zusammen.
Am Satzanfang steht \verb|\Cref|, im Satz \verb|\cref|.
Mehrere Ziele auf einmal: \cref{ch:einleitung,ch:fazit}.

\begin{table}[ht]
  \centering
  \begin{tabular}{lp{7.5cm}}
    \toprule
    Befehl & Wirkung \\
    \midrule
    \verb|\input{datei}|   & fügt die Datei an dieser Stelle ein \\
    \verb|\include{datei}| & wie \verb|\input|, zusätzlich neue Seite und eigene \texttt{.aux}-Datei \\
    \verb|\includeonly{...}| & übersetzt nur die genannten Dateien \\
    \bottomrule
  \end{tabular}
  \caption{Befehle zum Aufteilen eines Dokuments}
  \label{tab:befehle}
\end{table}

\section{Abkürzungen}

\gls{ci} kam schon in der Einleitung vor und erscheint hier nur noch als Kurzform.
\gls{cd} wird zum ersten Mal verwendet und daher ausgeschrieben.
\todo[inline]{Abschnitt zu Continuous Delivery ausbauen}

\chapter{Fazit}
\label{ch:fazit}

Wie in \cref{sec:verweise} gezeigt, erzeugt \verb|\cref| den passenden
Namen automatisch. Wird aus einer Tabelle eine Abbildung, ändert sich der Text
im Verweis von selbst.


% Abkürzungsverzeichnis ausgeben
\printnoidxglossary[type=\acronymtype,title=Abkürzungsverzeichnis]

\end{document}
